\begin{center}\small
\begin{tabular}{@{}lrrr@{}}
\toprule
Compact equation & Positive variables & Operations & Degree\\
\midrule
Package 142, raw parent &14&70&189\\
Package 142, reduced parent &14&71&157\\
Package 142, anchored parent &13&72&433\\
\midrule
New residue carrier $P_{13,70}$ &13&70&206\\
New residue carrier $P_{13,71}$ &13&71&162\\
New smaller-modulus $P_{13,72}$ &13&72&154\\
New rational carrier $P_{12,74}$ &12&74&438\\
\bottomrule
\end{tabular}
\end{center}
All seven rows in this table use the same complete indexed943/radix67
constant tuple. The first two new forms remove one coordinate without
adding an arithmetic gate. At thirteen variables and 72 gates, the
smaller-modulus equation reduces the matched anchored parent's degree
from 433 to 154. The twelve-variable form uses two additional gates
relative to that parent. These are exact circuit properties, not
estimates from the appearance of a displayed formula.

\begin{center}\small
\begin{tabular}{@{}lrrrr@{}}
\toprule
Direct integer quartic & Variables & Residuals & Monomials & Height\\
\midrule
Incoming Pooled Powers &225&162&7,867&128,446\\
Incoming Quadratic Projection &234&171&19,050&4,480\\
\midrule
New principal $P_{217}$ &217&154&19,112&5,184\\
New $P_{220}$, lower height &220&157&8,847&2,336\\
New $P_{220}$, smaller support &220&157&7,858&6,272\\
\bottomrule
\end{tabular}
\end{center}
The quartic heights are maximum absolute coefficients after complete
expansion and collection. They are not residual heights. All coordinates
are unrestricted integers, and all these polynomials have degree four.
The 217-variable result trades expanded support for fewer variables and
rows. The smaller-support 220-variable result improves all four displayed
size measures over its actual 225-variable parent. The lower-height
220-variable result makes a different support/height tradeoff.

The direct results are new compositions of existing arithmetic and
projection methods with the specified source changes. In particular the
incoming Linear CRT archive already uses $F=t^t$, $R=F$, and a linear
odd beta multiplier. We do not claim priority for those choices.
